\section{Data and experimental design}
\label{sec:data}

We use the EUPPBench postprocessing benchmark, version~\EuppbenchVersion{}
\citep{demaeyer_2023_euppbench}, restricted to its German subset of
\NStations{} stations and the \(2\,\mathrm{m}\) temperature target. EUPPBench
defines its own train/test split, which we never re-partition: training
uses ECMWF reforecasts issued \TrainYearRange{} with
\TrainEnsembleMembers{} ensemble members, and testing uses operational
ECMWF forecasts issued \TestYearRange{} with \TestEnsembleMembers{}
members. Both sides carry \NLeadTimes{} lead times from 0 to
\MaxLeadHours{}~h in \LeadStepHours{}-hour steps. Reference-quality
observations are DWD station measurements; the ensemble side is ECMWF's
operational and reforecast archives as distributed through EUPPBench.
EUPPBench also defines a wind-speed target (\(w10\)); we defer it here
because no country in the benchmark provides a direct wind-speed
observation series, only wind gust, and substituting gust for speed would
change what is being verified.

TimesFM-3's fixed context-window requirement (at least 40 prior
observations) excludes some station/lead/issue-time combinations that would
otherwise be evaluable by the other methods, so every comparison in this
paper is computed on the intersection of instances all methods can produce
a prediction for, not each method's own largest possible set. The full test
set otherwise available to EMOS and the raw ensemble contains 748{,}099
instances; requiring TimesFM-3's context window reduces this to
\NMatchedInstances{} matched instances, and every result in this paper is
computed on that matched set for every method, including EMOS and the raw
ensemble, not their own larger native set.

The archive's raw station dimension carries 51 station identifiers; two are
excluded throughout this study because their observation series contain
missing values that would otherwise propagate NaNs into every downstream
metric, leaving the \NStations{} stations used everywhere in this paper.
Ensemble summary statistics (mean and variance) are computed once from the
raw per-member forecasts before any postprocessing method sees them, using
an unbiased (\(\mathrm{ddof}=1\)) variance estimator on both the training
and test sides so that EMOS is never trained against a systematically
different variance definition than the one it is evaluated against. The
training-side ensemble variance is nonetheless estimated from
\TrainEnsembleMembers{} members against \TestEnsembleMembers{} on the test
side; an 11-member variance estimate is intrinsically noisier than a
51-member one of the same true spread, so EMOS is trained against a
noisier estimate of ensemble variance than the one used to evaluate it
(Section~\ref{sec:limitations}).

The paper's central axis is training-data size, expressed as a number of
training \emph{cases} \(k\) rather than calendar days directly, since the
reforecast archive samples a fixed, non-daily grid of issue dates rather
than every calendar day. The reforecast archive's raw \texttt{time} coordinate cannot be reliably
decoded to calendar dates: attempting to decode it raises an
\texttt{OverflowError}, confirmed to be caused by an unmasked fill-value
sentinel surviving into the time-unit conversion, not by a transient
network fault, since the same failure text was independently confirmed on
the sibling reforecast-observations file. We therefore treat each issue
date as an abstract
\((\mathrm{year\_idx}, \mathrm{time\_idx})\) pair rather than a decoded
calendar date: \ReforecastTotalIssueDates{} representative issue-date
positions (\(\mathrm{time\_idx}\)) recur once per analog year
(\(\mathrm{year\_idx}\), 20 analog years in total), so the full archive
contains \FullArchiveCases{} cases (\(\ReforecastTotalIssueDates \times
20\)). We measured the calendar-day span of one such 209-date cycle
directly from its real, decodable date labels: it spans
\ReforecastIssueDateSpanDays{} calendar days --
consistent with a twice-weekly issue cadence over two calendar years, not
one -- giving \(\DaysPerCase\) calendar days per case. We convert a
requested training size of \(N\) days to \(k = \mathrm{round}(N /
\DaysPerCase)\) cases throughout, and report both units together wherever
a training size is cited, since case count and calendar-day equivalent are
not interchangeable at this archive's actual sampling density. The
contiguous training arm (below) orders cases by descending
\((\mathrm{year\_idx}, \mathrm{time\_idx})\), assuming ascending
\(\mathrm{year\_idx}\) is chronologically more recent; only
\(\mathrm{time\_idx}\)'s ordering was independently verified against real
decoded dates, so this assumption is disclosed, not confirmed
(Section~\ref{sec:limitations} bounds its limited practical impact).

We test two ways of selecting which \(k\) cases make up a given training
size, since they answer different questions. The \emph{contiguous} arm
(primary) takes the \(k\) cases closest in time to the test period,
working backward -- the realistic scenario for a newly deployed station
whose limited history is chronologically contiguous, so small \(k\) sees
substantially less seasonal variety than a full archive would. The
\emph{random} arm (secondary, \DataSizeSweepSeeds{} seeds, mean and
standard deviation reported) draws \(k\) cases uniformly at random from
the full archive, giving artificial full-season coverage even at small
\(k\). The gap between the two arms' results is itself a reported finding
about the operational value of seasonal training coverage
(Section~\ref{sec:results}), not a discrepancy to be resolved in favor of
one arm. \citet{baran_2026_sharpness_penalty} targets the same sharpness
axis on the same benchmark; we discuss the relationship to our own
sharpness results directly in Section~\ref{sec:discussion} rather than
here. Other EUPPBench applications are surveyed in
Section~\ref{sec:intro}.
